\originalchapter{热大爆炸的困难、空宇宙与暴胀习题}
本章完成8.286 PS9七题与PS10四题。其中PS9.5和PS10.1题面相同，合并成一个习题并保留两个原编号；其余条件和子问均保留。解答为独立补写，数量级模型与实际数值积分区别说明。

\src{8.286 PS9.1，大爆炸核合成参数，(a)--(d)。}
\begin{exercise}
定性说明以下变化使预测氦丰度怎样改变，并给物理原因：
(a)增加一种中微子；
(b)弱反应更强，使$n/p$热平衡一直维持到$k_BT\simeq0.25$ MeV；
(c)质子中子质量差比1.29 $\mathrm{MeV}/c^2$更大；
(d)设早时相变将重子集中在小团块，中子能扩散至近乎均匀而带电质子多数留在团块，判断氦丰度的趋势。
\end{exercise}
\begin{solution}
简单估计的链条是：弱反应冻结留下中子比例，中子在氘瓶颈解除前部分衰变，余下中子多进入$^4$He，$Y_{\rm He}\simeq2(n/p)/(1+n/p)$。这些比较均保持其他参数不变。

(a)轻中微子提高固定$T$下的$g_*$和$H$，弱反应相对膨胀更早冻结，冻结温度更高、$n/p\simeq e^{-\Delta m c^2/(k_BT)}$更大；较短的等待时间还减少中子衰变，故氦增多。
(b)到更低温仍保持平衡使中子更少，$e^{-1.29/.25}\simeq0.0057$，所以氦显著下降。
(c)更大质量差更不利于中子，平衡比例下降，氦减少；完整网络还会考虑反应率的相应变化，此处只比较指定热力学效应。

(d)按题面中子、质子分离的图景，富质子团块缺少中子，均匀低密区中的中子又缺乏高密度质子反应伙伴，会有更多中子来不及结合而衰变，预计氦降低。这个结论依赖“分离保持到核反应阶段”的简化；实际不均匀核合成需同时求扩散、反扩散、团块尺度和反应网络，单凭非均匀性本身并不构成严格的单调定理。没有把这种趋势估计称为已执行的完整BBN网络。
\end{solution}

\src{8.286 PS9.2，电磁零点能的截止估计，(a)--(c)。}
\begin{exercise}
边长$L$周期盒内，电磁场只取$|\boldsymbol k|\le k_{\max}=2\pi/\lambda_{\min}$。
(a)在$\lambda_{\min}\ll L$时数模式；
(b)每模式最低能为$\hbar\omega/2$，取$\lambda_{\min}$等于Planck长度，求零点能的质量密度；
(c)与宇宙临界密度比较。周期条件$k_iL=2\pi n_i$，每$\boldsymbol k$有两个偏振。
\end{exercise}
\begin{solution}
每$k$空间格占$(2\pi/L)^3$，乘球体积及两个偏振，
\[
 N\simeq2\frac{L^3}{(2\pi)^3}\frac{4\pi k_{\max}^3}{3}
 =\frac{L^3k_{\max}^3}{3\pi^2}.
\]
正负波数分别为不同传播方向；不能先取一半再额外错去偏振因子。除体积后壳模式密度$\dd N/V=k^2\dd k/\pi^2$，积分零点能，
\[
 u_0=\int_0^{k_{\max}}\frac{\hbar ck}{2}\frac{k^2\dd k}{\pi^2}
 =\frac{\hbar c k_{\max}^4}{8\pi^2},\qquad
 \rho_0=\frac{\hbar k_{\max}^4}{8\pi^2c}.
\]
取$\ell_P=\sqrt{\hbar G/c^3}$、题面的$k_{\max}=2\pi/\ell_P$，得$\rho_0=2\pi^2c^5/(\hbar G^2)$。实际SI值$1.0175\times10^{98}\ \mathrm{kg\,m^{-3}}$；额外取$H_0=67.3$的今天临界密度，比例$1.1960\times10^{124}$。常见的$k_{\max}=1/\ell_P$是另一个截止约定，系数差$(2\pi)^4$，不能在本题中混用。

这是裸模式和的数量级困难：四次方UV敏感远大于观测宇宙密度。但它不是已求得可观测的宇宙常数。三动量硬截止本身不保持Lorentz不变，真实真空应讨论重整化、各场贡献和引力耦合；不能把这个数当作无需模型的精确暗能量预测。
\end{solution}

\src{8.286 PS9.3，1秒时的视界困难。}
\begin{exercise}
轻元素丰度成功暗示$t\simeq1$ s时已有热平衡。后来成为今天可观测宇宙的区域，在当时约跨多少个粒子视界？可假设空间平坦，给明确数量级估计。
\end{exercise}
\begin{solution}
必须区分“今天区域的半径”与“当时因果尺度”。为具体估算，额外采用PS10的今天有效可观测半径$R_0\simeq1.417\times10^4$ Mpc、$T_{\gamma0}=2.725$ K。恒定$g_*=10.75$辐射近似给1 s时$k_BT\simeq0.8592$ MeV。共动熵守恒，
\[
 a_1=\left(\frac{g_{*s,0}}{10.75}\right)^{1/3}
 \frac{k_BT_{\gamma0}}{0.8592\ {\rm MeV}},\quad
 g_{*s,0}=2+\frac78\,6\left(\frac4{11}\right)=3.9091,
\]
得到$a_1=1.9508\times10^{-10}$。该区域当时半径$R_0a_1=8.53\times10^{16}$ m，而辐射粒子视界$\ell_{p,h}=2ct=6.00\times10^8$ m。因此半径是$1.42\times10^8$个视界长度，直径跨$2.84\times10^8$个视界长度；按体积比较约有$10^{24}$量级的视界体积。

这不是说每个这样的区域都曾互相热化，恰恰指出常规辐射历史从$t=0$算起无法靠光速内的热过程建立整个区域的统一温度。用今天Hubble半径而非粒子视界半径会改变几倍系数，但不改变亿级跨越的困难。估计包含电子湮灭的熵效应，未假称没有粒子阈值的$aT$全程常数。
\end{solution}

\src{8.286 PS9.4，平坦性倒推。}
\begin{exercise}
假设今天$\Omega_0=0.1$，写早时$\Omega=1-\delta$，估$t=10^{-37}$ s时$\delta$，采用常规、未暴胀宇宙历史。
\end{exercise}
\begin{solution}
核心关系$1-\Omega=-kc^2/(a^2H^2)$；对开模型用今天$\Omega_{k0}>0$，
\[
 \delta=\frac{\Omega_{k0}H_0^2}{a^2H^2}.
\]
辐射主导$H=1/(2t)$，由$\dot a=H_0\sqrt{\Omega_{r0}}/a$积分$a^2=2H_0\sqrt{\Omega_{r0}}\,t$，故
\[
 \delta(t)\simeq\frac{2\Omega_{k0}H_0t}{\sqrt{\Omega_{r0}}}.
\]
为数值唯一，额外明确取$H_0=67.3$、$\Omega_{r0}=9.2\times10^{-5}$，不含真空，总$\Omega_0=0.1$意味着$\Omega_{m0}=0.1-\Omega_{r0}$、$\Omega_{k0}=0.9$。实际$\delta\simeq4.09\times10^{-53}$。若换今天归一、早时$g_*$与阈值，系数会变，$10^{-50}$余位的精细要求仍在。

原因是减速膨胀时$aH$随时间下降，微小曲率相对物质的重要性不断增长；不能仅凭早时Friedmann密度很大就说平坦性已被动态解释。暴胀的加速阶段反而使$aH$增长，曲率相对项被压低。
\end{solution}

\src{8.286 PS9.5与PS10.1，完全相同的单极题，(a)--(d)。}
\begin{exercise}
某GUT在$k_BT_c=10^{16}$ GeV相变，单极质量$M_Mc^2=10^{18}$ GeV。相变前平坦、辐射主导，$g_{\rm GUT}=200$，$\rho=(\pi^2g_{\rm GUT}/30)(k_BT)^4/(\hbar^3c^5)$。Kibble估$n_M\sim\xi^{-3}$，取相关长度$\xi$为物理粒子视界。
每问先给符号式再数值：(a)相变时间；(b)单极数密度；(c)$n_M/n_\gamma$；
(d)假设单极对湮灭可忽略，且该数目比一直与今天相同，估今天$\Omega_M$。只对产生这些单极的指定GUT模型作估计。
\end{exercise}
\begin{solution}
将题给密度代辐射年龄，
\[
 t_G=\sqrt{\frac{45}{16\pi^3g_{\rm GUT}}}
 \frac{\hbar^{3/2}c^{5/2}}{\sqrt G(k_BT_c)^2},\quad
 \xi=2ct_G,\quad n_M=(2ct_G)^{-3}.
\]
光子数密度$n_\gamma=[2\zeta(3)/\pi^2](k_BT_c/(\hbar c))^3$，因此
\[
 \frac{n_M}{n_\gamma}
 =\frac{\pi^2}{2\zeta(3)}
 (2ct_G)^{-3}\left(\frac{\hbar c}{k_BT_c}\right)^3,\quad
 \Omega_M=\frac{M_M\,n_{\gamma0}}{\rho_{c0}}\frac{n_M}{n_\gamma}.
\]
额外明确今天$T_{\gamma0}=2.725$ K、$H_0=67.3$。实际结果
\[
 t_G=1.7114\times10^{-39}\ {\rm s},\quad
 \xi=1.0261\times10^{-30}\ {\rm m},\quad
 n_M=9.2555\times10^{89}\ {\rm m^{-3}},
\]
\[
 n_\gamma=3.1703\times10^{94}\ {\rm m^{-3}},\quad
 n_M/n_\gamma=2.9195\times10^{-5},\quad
 \Omega_M\simeq2.51\times10^{21}.
\]
原题明确要求比值不变，所以上述(d)遵从它。真实光子会被相继消失的热物种加热，比值并非严格不变。若改保守的单极共动数和总熵，今天比值需乘$g_{*s,0}/g_{{*s},G}=3.9091/200$，得到$\Omega_M\simeq4.91\times10^{19}$；依然灾难性过大。这两个数对应不同假设，不能把熵修正后答案假称来自“不变比值”。

$\xi$只受因果上限约束，更小相关尺度会产生更多单极；$n_M$是数量级而非精确初值。暴胀可稀释既有单极，但再热温度过高会重新产生，且并非所有破缺模型都必须有相同拓扑缺陷。
\end{solution}

\src{8.286 PS9.6，空宇宙的FLRW描述，(a)--(d)，加分题。}
\begin{exercise}
令$\rho=0$、取$k=-1$，FLRW度量
$\dd s^2=-c^2\dd t^2+a^2(t)[\dd r^2/(1+r^2)+r^2\dd\Omega^2]$。
(a)求增长尺度因子；
(b)共动源$t_e$发射、$t_o$接收求红移；
(c)原点$t_e$发出的径向光求$r(t)$；
(d)以接收者$r$表示红移，说明为何与$t_e$无关。
\end{exercise}
\begin{solution}
约束$\dot a^2=c^2$，选坐标起点$a(0)=0$和增长支，$a=ct$。因而$1+z=t_o/t_e$。零测地线给
\[
 \frac{\dd r}{\sqrt{1+r^2}}=\frac{\dd t}{t},
 \qquad r(t)=\sinh\ln(t/t_e).
\]
接收时$\ln(t_o/t_e)=\operatorname{arsinh}r$，所以
\[
 1+z=e^{\operatorname{arsinh}r}=r+\sqrt{1+r^2}.
\]
时间整体缩放不改变两个时刻比例，因此同一共动$r$的红移不依赖发射时刻。此模型的粒子视界$ct\int_0^t\dd t'/t'$发散，与尘埃、辐射有限视界不同；$a=0$在此是坐标退化，不是有物质大爆炸的曲率奇点。
\end{solution}

\src{8.286 PS9.7，同一空宇宙的SR描述，(a)--(d)，加分题。}
\begin{exercise}
惯性球坐标$(t',r',\theta',\phi')$中，各共动观察者从原点沿恒速直线出发，限制在未来光锥内。宇宙时$t$为各观察者从原点开始的固有时。
(a)求$t(t',r')$；
(b)假定角坐标相同，以横向线段长度求$r(t',r')$，画等$t$与等$r$线；
(c)求$r$与惯性速度$v$关系，用SR红移验证上一题；
(d)写反变换及微分，代Minkowski间隔验证FLRW度量。
\end{exercise}
\begin{solution}
经过事件$(t',r')$的观察者$v=r'/t'$，故$t=t'\sqrt{1-v^2/c^2}=\sqrt{t'^2-r'^2/c^2}$。横向线段惯性长$r'\dd\theta'$等于FLRW长$ctr\dd\theta$，所以$r=r'/(ct)$。反变换为
\[
 t'=t\sqrt{1+r^2},\quad r'=ctr,\quad\theta'=\theta,\quad\phi'=\phi .
\]
等$t$是光锥内的双曲线$t'^2-r'^2/c^2=t^2$；等$r$是从原点发出的直线$r'/ct'=r/\sqrt{1+r^2}$：
\begin{center}\begin{tikzpicture}[scale=1.6]
\draw[->] (0,0)--(2.5,0) node[right] {$r'/c$};
\draw[->] (0,0)--(0,2.5) node[above] {$t'$};
\draw[gray,dashed] (0,0)--(2.3,2.3);
\foreach \b in {.25,.5,.75}{\draw[structurecolor] (0,0)--({2.3*\b},2.3);}
\foreach \a in {.45,.9,1.35}{\draw[black,domain=0:1.7,samples=70] plot(\x,{sqrt(\a*\a+\x*\x)});}
\end{tikzpicture}\end{center}
令$\beta=v/c$，$r=\beta/\sqrt{1-\beta^2}$。SR纵向退行红移$1+z=\sqrt{(1+\beta)/(1-\beta)}=\sqrt{1+r^2}+r$，与上一题一致。

再取微分
\[
 \dd t'=\sqrt{1+r^2}\dd t+\frac{tr}{\sqrt{1+r^2}}\dd r,\qquad
 \dd r'=cr\dd t+ct\dd r.
\]
代$-c^2\dd t'^2+\dd r'^2$，交叉项恰取消，得到$-c^2\dd t^2+c^2t^2\dd r^2/(1+r^2)$；角项$r'^2\dd\Omega'^2=c^2t^2r^2\dd\Omega^2$。故空开FLRW只是Minkowski未来光锥的一套坐标。任何非零尘埃／辐射早时都会使密度上升，不能因今天很稀就把此空模型延伸到真实宇宙的最早时期。
\end{solution}

\src{8.286 PS10.2，任意曲率的真空增长解。}
\begin{exercise}
恒定假真空密度$\rho_f$，$\chi=\sqrt{8\pi G\rho_f/3}$，满足$(\dot a/a)^2=\chi^2-kc^2/a^2$。求任意$k$的增长解，分别考虑正负曲率，并证明足够晚时均趋$a\propto e^{\chi t}$。
\end{exercise}
\begin{solution}
等价$\dot a^2=\chi^2a^2-kc^2$。对$k>0$最小半径$c\sqrt{k}/\chi$，增长支从反弹时刻$t_b$开始：
\[
 a=\frac{c\sqrt{k}}{\chi}\cosh[\chi(t-t_b)],\qquad t>t_b .
\]
对$k<0$可选$a=0$坐标时刻$t_s$，
\[
 a=\frac{c\sqrt{-k}}{\chi}\sinh[\chi(t-t_s)],\qquad t>t_s .
\]
对$k=0$，$a=Ae^{\chi t}$。逐式微分、平方可验证同一约束；正曲率的$\cosh$不是从$a=0$爆炸支增长。因$\sinh x,\cosh x\sim e^x/2$，三者晚时都指数增长。更一般$a=Ae^{\chi t}+Be^{-\chi t}$需$AB=kc^2/(4\chi^2)$，取正尺度因子与相应增长区间。这个结论仅在假真空密度持续恒定期间，不是已经解释退出和再热。
\end{solution}

\src{8.286 PS10.3，今天的有效粒子视界，(a)--(b)。}
\begin{exercise}
(a)用闭空间度量$-c^2\dd t^2+\tilde a^2(t)[\dd\psi^2+\sin^2\psi\dd\Omega^2]$，$\tilde a=a/\sqrt{k}$，证明光路视界$R_{p,h}=a(t)\int_0^tc\dd t'/a(t')$对曲率仍成立；
(b)以2013题面参数$H_0=67.3,\Omega_{m0}=.315,\Omega_{v0}=.685,\Omega_{r0}=9.2\times10^{-5},T_{\gamma0}=2.725$ K实际求今天视界，给Mpc和光年。积分不含暴胀阶段，指暴胀后可观测光的有效视界。
\end{exercise}
\begin{solution}
径向零测地线$\dd\psi=c\dd t/\tilde a$，今天沿它积累的实际路长为$\tilde a_0\psi_h$；$\sqrt{k}$在两因子中取消，即给题式。闭空间路径若绕行，积累路长不等于最短两点距离，前者才是这个积分的定义。

取$x=a/a_0$并用连续方程，
\[
 R_{p,h,0}=\frac c{H_0}\int_0^1
 \frac{\dd x}{\sqrt{\Omega_{r0}+\Omega_{m0}x+\Omega_{k0}x^2+\Omega_{v0}x^4}}.
\]
题面小数之和为1.000092；直接保留它们并强制Friedmann闭合给$\Omega_{k0}=-9.2\times10^{-5}$。实际积分3.180072696，故
\[
 R_{p,h,0}=14165.85\ {\rm Mpc}=4.62028\times10^{10}\ {\rm ly}.
\]
若把题意视为平坦且把真空项调整为$1-\Omega_m-\Omega_r$，结果14165.54 Mpc，差约0.0022\%；这是参数舍入处理的区别，不是声称Planck测定了这几十万分之一的精度。相同背景的年龄约13.812 Gyr，视界光行长度并不等于$ct_0$。

若把真实早时暴胀也纳入原始粒子视界，积分起点与尺度因子历史会改变；这里按题意只算普通热大爆炸后光的有效可观测半径。
\end{solution}

\src{8.286 PS10.4，解决均匀性所需的膨胀，加分题。}
\begin{exercise}
初$t_i$近乎均匀、平坦，假真空$\rho_f=E_f^4/(\hbar^3c^5)$，$E_f=E_{16}10^{16}$ GeV，$H_i=\sqrt{8\pi G\rho_f/3}$。初已平滑长度$r_h(t_i)=\beta c/H_i$、$0<\beta\lesssim1$。暴胀使$a$增长$Z$，瞬间结束并再热，全部能量变为$g_{\rm RH}$种同温相对论辐射。此后只随尺度因子拉伸均匀区、共动熵守恒。以题面PS10.3今天视界，求$r_h(t_0)>R_{p,h,0}$所需$Z_{\min}(E_{16},g_{\rm RH},\beta)$及$N_{\min}=\ln Z_{\min}$；今天熵包括较冷中微子。
\end{exercise}
\begin{solution}
能量匹配给$k_BT_{\rm RH}=(30/(\pi^2g_{\rm RH}))^{1/4}E_f$。采用再热时能量与熵自由度相等这一题设近似，熵守恒为
\[
 \frac{a_e}{a_0}
 =\left(\frac{g_{*s,0}}{g_{\rm RH}}\right)^{1/3}
 \frac{T_{\gamma0}}{T_{\rm RH}},\qquad
 g_{*s,0}=2+\frac78\,6(T_\nu/T_\gamma)^3=\frac{43}{11}.
\]
今天均匀长度$r_h(t_0)=(a_0/a_e)Z\beta c/H_i$，所以
\[
 Z_{\min}=\frac{R_{p,h,0}H_i}{\beta c}
 \left(\frac{g_{*s,0}}{g_{\rm RH}}\right)^{1/3}
 \frac{k_BT_{\gamma0}}{(30/(\pi^2g_{\rm RH}))^{1/4}E_f}.
\]
因$H_i\propto E_f^2$，比例为$E_f g_{\rm RH}^{-1/12}\beta^{-1}$，不是$E_f^2$或很强的粒子数依赖。实际常数代入：
\[
 Z_{\min}=9.1465\times10^{26}
 \frac{E_{16}}{\beta}\left(\frac{g_{\rm RH}}{300}\right)^{-1/12},
\]
\[
 N_{\min}=62.0806+\ln E_{16}-\ln\beta
 -\frac1{12}\ln(g_{\rm RH}/300).
\]
严格的大于号要求实际$N>N_{\min}$。该例$H_i^{-1}=2.7764\times10^{-38}$ s、再热$k_BT_{\rm RH}=3.173\times10^{15}$ GeV，均由同一预算计算。这个$N$是把指定初始均匀区拉到今天视界所需的总膨胀；它与某个观测模式离开Hubble尺度至退出的e折数、永恒暴胀持续时间不是同一问题。初$\beta$、再热持续时间及额外熵产生会改变门槛。
\end{solution}
